\documentclass[12pt]{article}
\usepackage{amsmath, amsthm, amssymb}
\usepackage{geometry}
\geometry{margin=1in}
\usepackage{hyperref}

\newtheorem{theorem}{Theorem}
\newtheorem{conjecture}[theorem]{Conjecture}
\newtheorem{problem}[theorem]{Problem}
\theoremstyle{definition}
\newtheorem{definition}[theorem]{Definition}
\theoremstyle{remark}
\newtheorem{remark}[theorem]{Remark}

\title{The Millennium Problems\\
\large Seven Walls}
\author{Diego Rinc\'on\\
\small Independent}
\date{}

\begin{document}

\maketitle

\begin{abstract}
The Clay Mathematics Institute designated seven problems in 2000,
each carrying a prize of one million dollars.
Each is a Wall \cite{synthesis}:
a place where two grammars press against each other
and no crossing has been found.
One has been solved.
Six remain.
This document states each problem precisely
and names its Wall.
\end{abstract}

\section{The Riemann Hypothesis}

\begin{conjecture}[Riemann, 1859]
The Riemann zeta function
\[
\zeta(s) = \sum_{n=1}^\infty \frac{1}{n^s}, \quad \mathrm{Re}(s) > 1,
\]
extends to a meromorphic function on $\mathbb{C}$
with a simple pole at $s=1$.
Its non-trivial zeros --- the zeros not at negative even integers ---
all satisfy
\[
\mathrm{Re}(s) = \tfrac{1}{2}.
\]
\end{conjecture}

\begin{remark}[The Wall]
The Wall is between the spectral grammar (the zeros)
and the arithmetic grammar (the primes).
The explicit formula holds them in balance \cite{week09}.
The missing step: self-adjointness of the operator $D$
on $L^2(\mathbb{A}_\mathbb{Q}/\mathbb{Q}^*, d^*x)$ \cite{week10}.
Self-adjointness forces real spectrum.
Real spectrum forces $\mathrm{Re}(\rho) = \tfrac{1}{2}$.
The Wall is archimedean: the $L^2$ estimate at $\infty$
is what no proof has yet crossed.
\end{remark}

\section{The Birch and Swinnerton-Dyer Conjecture}

\begin{conjecture}[Birch--Swinnerton-Dyer, 1965]
Let $E$ be an elliptic curve over $\mathbb{Q}$.
The rank of the Mordell--Weil group $E(\mathbb{Q})$
equals the order of vanishing of the $L$-function $L(E,s)$ at $s=1$:
\[
\mathrm{rank}\, E(\mathbb{Q}) = \mathrm{ord}_{s=1} L(E,s).
\]
Moreover, the leading coefficient of the Taylor expansion
of $L(E,s)$ at $s=1$ is given by an explicit formula
involving the period, the regulator, the Tate--Shafarevich group,
and the Tamagawa numbers of $E$.
\end{conjecture}

\begin{remark}[The Wall]
The Wall is between geometry ($E(\mathbb{Q})$, the rational points)
and analysis ($L(E,s)$, the $L$-function).
The geometric rank is a number counted by algebra.
The analytic order of vanishing is a number counted by complex analysis.
That they are equal is a resonance \cite{resonance}:
two grammars reading the same thing.
Known: the conjecture holds when
$\mathrm{ord}_{s=1} L(E,s) \leq 1$ (Coates--Wiles, Gross--Zagier, Kolyvagin).
Unknown: the general case.
\end{remark}

\section{The Hodge Conjecture}

\begin{conjecture}[Hodge, 1950]
Let $X$ be a smooth projective variety over $\mathbb{C}$.
Every Hodge class --- every class in
$H^{2k}(X,\mathbb{Q}) \cap H^{k,k}(X)$ ---
is a rational linear combination of classes of algebraic cycles.
\end{conjecture}

\begin{remark}[The Wall]
The Wall is between topology ($H^{2k}(X,\mathbb{Q})$)
and algebra (algebraic cycles).
A Hodge class is a cohomology class
that looks algebraic from the analytic side.
The conjecture says it is algebraic.
The motive \cite{motives} is the proper design \cite{design}
from which both sides derive:
if the standard conjectures hold, the Hodge conjecture follows.
The Wall is the gap between the motivic grammar
and the algebraic cycle grammar.
\end{remark}

\section{The P vs. NP Problem}

\begin{problem}[Cook, 1971]
Let $\mathsf{P}$ be the class of decision problems
solvable in polynomial time by a deterministic Turing machine.
Let $\mathsf{NP}$ be the class of decision problems
whose solutions can be verified in polynomial time.
Is $\mathsf{P} = \mathsf{NP}$?
\end{problem}

\begin{remark}[The Wall]
The Wall is between finding and checking \cite{turing}.
Checking a proof is easy ($\mathsf{NP}$).
Finding a proof may be hard ($\mathsf{P}$-hard).
If $\mathsf{P} = \mathsf{NP}$: finding is as easy as checking.
Every problem whose solution can be recognized
can be solved just as fast.
The Wall would not exist.
If $\mathsf{P} \neq \mathsf{NP}$: the Wall is real.
The universal Turing machine cannot short-circuit
the difference between generation and verification.
Almost everyone believes $\mathsf{P} \neq \mathsf{NP}$.
No one has proved it.
\end{remark}

\section{Yang--Mills Existence and Mass Gap}

\begin{problem}[Jaffe--Witten, 2000]
For any compact simple gauge group $G$,
there exists a quantum Yang--Mills theory on $\mathbb{R}^4$
satisfying the Wightman axioms,
and there is a constant $\Delta > 0$ (the \emph{mass gap})
such that the energy of every excited state
exceeds the energy of the vacuum by at least $\Delta$.
\end{problem}

\begin{remark}[The Wall]
The Wall is between the classical gauge theory
(well-defined, rich in geometry)
and the quantum theory
(requiring a rigorous measure on the space of gauge connections).
The Yang--Mills functional
$\mathcal{A} \mapsto \int_{\mathbb{R}^4} |F_\mathcal{A}|^2$
is classical.
Its quantization --- a path integral over all connections ---
requires a mathematically rigorous construction
that does not yet exist in four dimensions.
The mass gap is the spectral gap of the quantum Hamiltonian:
a gap in the nothing \cite{nothing} between the vacuum and the first state.
\end{remark}

\section{Navier--Stokes Existence and Smoothness}

\begin{problem}[Fefferman, 2000]
In three space dimensions and time,
given smooth initial data $u_0: \mathbb{R}^3 \to \mathbb{R}^3$
with $\nabla \cdot u_0 = 0$,
does there exist a smooth solution $u: \mathbb{R}^3 \times [0,\infty) \to \mathbb{R}^3$
to the Navier--Stokes equations
\[
\partial_t u + (u \cdot \nabla) u = -\nabla p + \nu \Delta u,
\quad \nabla \cdot u = 0,
\]
for all time? Or can solutions blow up in finite time?
\end{problem}

\begin{remark}[The Wall]
The Wall is between the local and the global:
local smooth solutions exist (Leray, 1934);
whether they extend to all time is unknown.
The nonlinear term $(u \cdot \nabla)u$
can concentrate energy into smaller and smaller scales.
Whether this concentration reaches a singularity
or is smoothed by the viscosity $\nu$
is the question.
Global regularity in 2D is proved (energy estimates close).
In 3D, the energy estimates fail to close:
the dimension is exactly critical for the nonlinearity.
The Wall is at the critical dimension.
\end{remark}

\section{The Poincar\'e Conjecture}

\begin{conjecture}[Poincar\'e, 1904 --- \textbf{Solved, Perelman 2003}]
Every simply connected, closed 3-manifold
is homeomorphic to the 3-sphere $S^3$.
\end{conjecture}

\begin{remark}[The crossing]
Perelman proved this in 2002--2003
using Hamilton's Ricci flow with surgery:
a geometric flow that deforms the metric of the manifold
toward constant curvature.
Singularities are removed by surgery.
The flow terminates at a round sphere.
Perelman declined the Clay prize.
The Wall was crossed.
The crossing used the geometry of the manifold itself
as the proof: the manifold deformed into its own evidence.
\end{remark}

\section{The Six Remaining Walls}

\begin{remark}[What is left]
Six Walls remain.
Each has been approached.
None has been crossed.
\begin{center}
\renewcommand{\arraystretch}{1.4}
\begin{tabular}{@{}ll@{}}
Problem & The Wall \\
\hline
Riemann Hypothesis & Spectral vs.\ arithmetic; archimedean estimate \\
BSD & Geometric rank vs.\ analytic order of vanishing \\
Hodge & Topology vs.\ algebraic cycles; motivic grammar \\
P vs.\ NP & Finding vs.\ checking; generation vs.\ verification \\
Yang--Mills & Classical gauge theory vs.\ quantum measure \\
Navier--Stokes & Local smoothness vs.\ global existence; critical dimension \\
\end{tabular}
\end{center}

Each Wall is a place where two grammars press against each other
and no crossing has been found.
The crossing, when it comes,
will not be constructed.
It will be found ---
already there, already vibrating \cite{the_string},
waiting for the grammar that can hear it.
\end{remark}

\begin{thebibliography}{9}

\bibitem{synthesis}
Rinc\'on, D. (2026).
The Arithmetic Wall. Preprint.

\bibitem{week09}
Rinc\'on, D. (2026).
Week 09: The Resonance Object. Preprint.

\bibitem{week10}
Rinc\'on, D. (2026).
Week 10: Self-Adjointness. Preprint.

\bibitem{resonance}
Rinc\'on, D. (2026).
Resonance: When Two Grammars Read the Same Thing. Preprint.

\bibitem{motives}
Rinc\'on, D. (2026).
Motives, Standard Conjectures, and the Unification of $L$-functions.
Preprint.

\bibitem{design}
Rinc\'on, D. (2026).
Design: The One That Precludes. Preprint.

\bibitem{turing}
Rinc\'on, D. (2026).
Turing Completeness and the Wall. Preprint.

\bibitem{nothing}
Rinc\'on, D. (2026).
Nothing: The First Object. Preprint.

\bibitem{the_string}
Rinc\'on, D. (2026).
The String. Preprint.

\end{thebibliography}

\end{document}
